\originalchapter{Gaussian关联与Ginzburg准则}
\label{ch:gaussian}
\source{8.334 L4--L5及L8的Gaussian直接解。}
\originalsection{Gaussian积分与生成泛函}
有限维正定矩阵$A$有$Z[J]=\int\dd^nx\exp[-x^TAx/2+J^Tx]=(2\pi)^{n/2}(\det A)^{-1/2}\exp[J^TA^{-1}J/2]$。证明是完成平方$x\to x-A^{-1}J$再以正交矩阵对角化；若有零本征值或负本征值，积分并不存在为同一有限数。连续场在截断动量或有限晶格上使用同一公式，最后再取受控极限。

对$\mathcal H_0=\frac12\int[K(\nabla\phi)^2+r\phi^2]$，$G(q)=\avg{\phi_q\phi_{-q}}=1/(Kq^2+r)$、$\xi=\sqrt{K/r}$。配分自由能密度的Gaussian涨落项$\frac12\int\log(Kq^2+r)\dd^dq/(2\pi)^d$包含截止依赖的常数和解析背景，物理临界奇异部分需与之区分。

Wick定理由$Z[J]$的指数二次形式直接生成：奇数零均值矩为0，偶数矩为所有配对的$G$乘积。四点$\avg{1234}=G_{12}G_{34}+G_{13}G_{24}+G_{14}G_{23}$；连通四点在纯Gaussian模型为0，加入相互作用后才出现。
\originalsection{实空间相关与散射}
三维Fourier反演$G(r)=\e^{-r/\xi}/(4\pi Kr)$，可通过$(-K\nabla^2+r)G=\delta$的径向Green方程得到。临界$r\to0$，一般$d>2$的长距离$G(x)\propto1/|x|^{d-2}$。因此Gaussian指数$\eta=0,\nu=1/2$，静态磁化率$\chi=G(q=0)=1/r$。这里$r$同时表示质量参数和有时实空间距离容易混淆，正文距离用$|x|$或明确下标。

实验小角散射测结构因子$S(q)$近$1/(q^2+\xi^{-2})$的Lorentz峰，临界相关长度增长使峰变窄。真实临界$S(q)\sim q^{-2+\eta}$，非零$\eta$表示长距离场尺度已被相互作用修正。有限分辨率、样品尺寸与非临界背景须先校正，再把峰宽解释为$\xi^{-1}$。
\originalsection{涨落修正的尺度}
有序相鞍点$\phi_0^2=|r|/u$。一个相关体积内粗粒化的涨落方差主要来自$q\lesssim\xi^{-1}$，$\avg{(\delta\phi)^2}_\xi\sim\int_0^{1/\xi}\dd^dq/[Kq^2+O(|r|)]\sim K^{-d/2}|r|^{(d-2)/2}$。将它与$\phi_0^2$比较，
\begin{equation}
\frac{\avg{(\delta\phi)^2}_\xi}{\phi_0^2}\sim uK^{-d/2}|r|^{(d-4)/2}\ll1
\end{equation}
是平均场自洽条件。$d>4$临界附近比值下降，$d<4$增长，$d=4$出现边际对数。这给上临界维数4。它不是简单比较某点的全紫外方差；后者通常含大的非临界截止项。

一回路$tadpole$修正$\delta r\propto u\int\dd^dq/(Kq^2+r)$可移临界温度，导数$\partial_r\delta r\propto-u\int\dd^dq/(Kq^2+r)^2$在$d\le4$有红外奇异。这与Ginzburg比较是同一失控尺度的两种表现。质量偏移依微观截止，普适临界指数由后续尺度流的线性化得到。
\originalsection{下与上临界维数}
下临界维数问涨落或缺陷是否连有序相都破坏，上临界维数问临界指数何时由Gaussian/平均场决定。连续$O(n>1)$的热自旋波下临界2，标量Ising下临界1，普通短程$\phi^4$上临界4。二维XY在边界维数仍可有代数序与BKT，不能把“无通常长程序”解释成所有相关都短程；二维Ising则在其下临界之上、上临界之下，有非Gaussian有限温临界点。
